% ---------------------------------------------------------------------
% Resumo (NBR 6028: parágrafo único, 150 a 500 palavras, voz ativa)
% ---------------------------------------------------------------------
\setlength{\absparsep}{18pt}
\begin{resumo}
Dados de monitoramento de chuva e vento usados pela Defesa Civil são dados críticos: um valor perdido, adulterado ou atrasado pode atrasar um alerta de deslizamento ou alagamento. Este trabalho apresenta o projeto e o protótipo de um banco de dados de telemetria IoT em PostgreSQL, aplicado a duas estações meteorológicas simuladas em São José dos Campos (SP), que medem temperatura e umidade do ar, chuva, direção e velocidade do vento e radiação solar a cada minuto e transmitem lotes a cada dez minutos. A partir de uma versão inicial focada em desempenho, o grupo levantou requisitos de dados críticos e implementou, no próprio banco, a validação de cada leitura com datas vindas somente do relógio do servidor, eventos imutáveis, papéis de mínimo privilégio, isolamento por cliente com \textit{Row Level Security}, senhas com bcrypt e bloqueio por tentativas, e uma trilha de auditoria somente de inclusão encadeada por hash SHA-256. Um servidor em Node.js e um painel em HTML, CSS e JavaScript puros completam o protótipo, com atualização em tempo real, backups diários verificados e uma fila de contingência que guarda os lotes quando o banco está fora do ar. Seis meses de operação (cerca de 529 mil leituras) foram simulados e os mecanismos foram testados: a adulteração proposital da auditoria foi detectada no registro exato, o PostgreSQL foi derrubado durante o envio de dados sem perda de leituras, e uma cópia de segurança foi restaurada com as mesmas contagens e a trilha de auditoria íntegra. As consultas mantiveram os ganhos das otimizações (até 200 vezes com view materializada), e as validações aumentaram em cerca de 30\,\% o tempo de envio de uma leitura por vez, com impacto desprezível em lotes de dez. Conclui-se que um banco relacional bem modelado pode concentrar as garantias de integridade, confidencialidade, rastreabilidade e disponibilidade exigidas por dados críticos, sem depender da boa conduta da aplicação.

\vspace{\onelineskip}
\noindent\textbf{Palavras-chave}: banco de dados; dados críticos; PostgreSQL; auditoria; IoT.
\end{resumo}

\begin{resumo}[Abstract]
\begin{otherlanguage*}{english}
Rain and wind monitoring data used by civil defense agencies are critical data: a lost, tampered or delayed value may delay a landslide or flood warning. This work presents the design and prototype of an IoT telemetry database in PostgreSQL, applied to two simulated weather stations in São José dos Campos, Brazil, which measure air temperature and humidity, rainfall, wind direction and speed, and solar radiation every minute and transmit batches every ten minutes. Starting from a performance-oriented first version, the group gathered critical-data requirements and implemented, inside the database, validation of every reading with timestamps taken only from the server clock, immutable events, least-privilege roles, per-customer isolation with Row Level Security, bcrypt passwords with lockout, and an append-only audit trail chained by SHA-256 hashes. A Node.js server and a plain HTML, CSS and JavaScript dashboard complete the prototype, with real-time updates, verified daily backups and a contingency queue that keeps batches while the database is down. Six months of operation (about 529 thousand readings) were simulated and the mechanisms were tested: deliberate tampering of the audit trail was detected at the exact record, PostgreSQL was shut down while data was being sent without losing readings, and a backup was restored with identical row counts and an intact audit trail. Queries kept the gains of the optimizations (up to 200 times with a materialized view), and validation increased single-reading inserts by about 30\,\%, with negligible impact on ten-reading batches.

\vspace{\onelineskip}
\noindent\textbf{Keywords}: database; critical data; PostgreSQL; audit trail; IoT.
\end{otherlanguage*}
\end{resumo}
